\documentclass[troisieme]{fiche}
\metadonnees{12}{Tout est vert}{Bloc C — Décrire et comparer}

\begin{document}
\entete

\objectifs{%
\textbf{Langue} : le superlatif ; \emph{one of the\ldots est}, \emph{ever}.\par
\textbf{Texte} : L.\,F. Baum, \emph{The Wonderful Wizard of Oz} (1900), ch. \textsc{xi}.\par
\textbf{Durée} : 1 h — exercices 20 min, texte et questions 25 min, version 15 min.}

%% ===================================================================
\exercice[12 min]{Le superlatif}

\rappel{\textbf{Rappel.} Adjectif court : \emph{the\ldots -est}. Adjectif long :
\emph{the most\ldots} Le \emph{the} ne tombe jamais.}

\consigne{Complétez avec le superlatif de l'adjectif entre parenthèses.}

\begin{anglais}
\begin{enumerate}[label=\arabic*.,itemsep=2.5mm,leftmargin=*]
  \item The Emerald City is \trou{the most beautiful} city in Oz. \emph{(beautiful)}
  \item Dorothy was \trou{the youngest} of the four travellers. \emph{(young)}
  \item It was \trou{the greenest} place she had ever seen. \emph{(green)}
  \item The Lion says he is \trou{the biggest} coward in the world. \emph{(big)}
  \item This is \trou{the best} lemonade in the street. \emph{(good)}
  \item Of all their plans, that one was \trou{the worst}. \emph{(bad)}
  \item Toto is \trou{the smallest} of them all. \emph{(small)}
  \item That was \trou{the most surprising} answer of the day. \emph{(surprising)}
  \item The old housekeeper wore \trou{the dullest} dress in the house. \emph{(dull)}
  \item Which of you is \trou{the happiest}? \emph{(happy)}
\end{enumerate}
\end{anglais}

\note[Former le superlatif]{Même partage qu'au comparatif : \emph{-est} pour les courts (2, 3,
4, 7, 10), \emph{the most} pour les longs (1, 8). Doublement de la consonne finale après
voyelle brève : \emph{big → biggest} (4). \emph{Good → the best}, \emph{bad → the worst}
(5, 6).}

\note[Le complément du superlatif]{Un lieu s'introduit par \emph{in} : \emph{the biggest
coward in the world}, \emph{the best lemonade in the street}. Un ensemble de personnes ou de
choses s'introduit par \emph{of} : \emph{the youngest of the four travellers}.}

%% ===================================================================
\exercice[8 min]{Superlatifs en contexte}

\rappel{\textbf{Rappel.} \emph{One of the} + superlatif + pluriel. Avec \emph{ever}, le verbe
se met au present perfect : \emph{the best book I have ever read}.}

\consigne{Traduisez en anglais.}

\begin{anglais}
\begin{enumerate}[label=\arabic*.,itemsep=3mm,leftmargin=*]
  \item C'est le jour le plus long de l'année.
    \lignes{1}
    \corr{It is the longest day of the year.}
  \item C'est l'une des plus vieilles maisons de la ville.
    \lignes{1}
    \corr{It is one of the oldest houses in the town.}
  \item C'est le meilleur livre que j'aie jamais lu.
    \lignes{1}
    \corr{It is the best book I have ever read.}
  \item Le plus difficile est de commencer.
    \lignes{1}
    \corr{The hardest thing is to begin.}
  \item C'est la deuxième ville du pays.
    \lignes{1}
    \corr{It is the second largest city in the country.}
  \item C'est la plus jeune de la famille.
    \lignes{1}
    \corr{She is the youngest in the family.}
\end{enumerate}
\end{anglais}

\note[Trois pièges]{\emph{One of the oldest houses} : pluriel obligatoire après \emph{one of
the}. \emph{The best book I have ever read} : pas de pronom relatif, et present perfect avec
\emph{ever}. \og La deuxième ville \fg{} se dit \emph{the second largest city} : l'anglais
nomme l'adjectif que le français sous-entend.}

\note[\og Le plus difficile \fg]{Le français substantive l'adjectif ; l'anglais ne le fait pas
et ajoute un nom : \emph{the hardest thing}, \emph{the best part}. \emph{The most difficult is
to begin} n'est pas correct.}

%% ===================================================================
\newpage
\section{Texte}

\noindent\small\textit{Dorothy, l'Épouvantail, le Bûcheron de fer-blanc et le Lion arrivent aux
portes de la Cité d'Émeraude. Le gardien leur a fait chausser des lunettes vertes, fermées par
une petite clef, avant de les laisser entrer.}\normalsize

\begin{texte}
Even with eyes protected by the green \gl[spectacles]{spectacles}{des lunettes}, Dorothy and
her friends were at first \gl[to dazzle]{dazzled}{éblouir} by the
\gl[brilliancy]{brilliancy}{l'éclat} of the wonderful City. The streets were
\gl[to be lined with]{lined with}{être bordé de} beautiful houses all built of green
\gl[marble]{marble}{le marbre} and \gl[to stud]{studded}{parsemer, clouter} everywhere with
sparkling emeralds. They walked over a \gl[a pavement]{pavement}{un pavage, un trottoir} of
the same green marble, and where the blocks were joined together were rows of emeralds, set
closely, and \gl[to glitter]{glittering}{scintiller} in the brightness of the sun. The window
\gl[a pane]{panes}{une vitre} were of green glass; even the sky above the City had a green
\gl[a tint]{tint}{une teinte}, and the rays of the sun were green.

There were many people --- men, women, and children --- walking about, and these were all
dressed in green clothes and had greenish skins. They looked at Dorothy and her strangely
\gl[assorted]{assorted}{assorti, dépareillé} company with wondering eyes, and the children all
ran away and hid behind their mothers when they saw the Lion; but no one spoke to them. Many
shops stood in the street, and Dorothy saw that everything in them was green. Green candy and
green pop-corn were offered for sale, as well as green shoes, green hats, and green clothes of
all sorts. At one place a man was selling green lemonade, and when the children bought it
Dorothy could see that they paid for it with green \gl[a penny]{pennies}{un sou}.

There seemed to be no horses nor animals of any kind; the men carried things around in little
green \gl[a cart]{carts}{une charrette}, which they pushed before them. Everyone seemed happy
and contented and \gl[prosperous]{prosperous}{prospère}.

The Guardian of the Gates led them through the streets until they came to a big building,
exactly in the middle of the City, which was the Palace of Oz, the Great Wizard. There was a
soldier before the door, dressed in a green uniform and wearing a long green
\gl[a beard]{beard}{une barbe}.

``Here are strangers,'' said the Guardian of the Gates to him, ``and they demand to see the
Great Oz.''

``Step inside,'' answered the soldier, ``and I will carry your message to him.''
\end{texte}
\source{L.\,F. Baum, \emph{The Wonderful Wizard of Oz}, Chicago, Hill, 1900, ch.~\textsc{xi}.}

\partiel[15 min]{Questions}
\consigne{Répondez aux deux premières questions en français, aux deux dernières en anglais.}

\begin{enumerate}[label=\arabic*.,itemsep=2.5mm,leftmargin=*]
  \item Relevez cinq choses qui sont vertes dans la ville.
    \lignes{3}
    \corr{Les maisons de marbre vert, le pavage, les vitres, le ciel, les rayons du soleil,
    les vêtements et la peau des habitants, les bonbons, le pop-corn, les chaussures, les
    chapeaux, la limonade, les sous, les charrettes, l'uniforme et la barbe du soldat.}
  \item Comment les habitants accueillent-ils les voyageurs ?
    \lignes{2}
    \corr{Ils les regardent avec des yeux étonnés, sans leur adresser la parole ; les enfants
    s'enfuient et se cachent derrière leur mère en voyant le Lion.}
  \item \en{The travellers wear green spectacles, locked on with a key. What does that detail
    suggest about the city?}
    \lignes{3}
    \corr{That the green may be in the glasses and not in the city. The narrator never says
    the city is green: he says what Dorothy sees through the spectacles she cannot take off.
    The reader is told the trick long before Dorothy understands it.}
  \item \en{The verb \emph{seem} appears twice in the third paragraph. What does it add?}
    \lignes{2}
    \corr{\emph{There seemed to be no horses} ; \emph{Everyone seemed happy}. Le narrateur ne
    garantit rien : il rapporte une apparence. Le mot prépare la découverte que toute la cité
    est une illusion.}
\end{enumerate}

%% ===================================================================
\section{Version}
\consigne{Traduisez le passage en français. Il précède immédiatement le texte : le gardien des
portes prépare les voyageurs.}

\begin{version}
He opened the big box, and Dorothy saw that it was filled with spectacles of every size and
shape. All of them had green glasses in them. The Guardian of the Gates found a pair that
would just fit Dorothy and put them over her eyes. There were two golden
\gl[a band]{bands}{une bande, un bandeau} fastened to them that passed around the back of her
head, where they were locked together by a little key that was at the end of a
\gl[a chain]{chain}{une chaîne} the Guardian of the Gates wore around his neck. When they were
on, Dorothy could not take them off had she wished, but of course she did not wish to be
blinded by the \gl[a glare]{glare}{un éclat aveuglant} of the Emerald City, so she said
nothing.
\end{version}
\source{\emph{Ibid.}, ch.~\textsc{x}.}

\lignes{14}

\corr{\textbf{Proposition de traduction.} Il ouvrit la grande boîte, et Dorothy vit qu'elle
était pleine de lunettes de toutes les tailles et de toutes les formes. Toutes avaient des
verres verts. Le Gardien des Portes trouva une paire qui allait exactement à Dorothy et la lui
posa sur les yeux. Deux bandeaux d'or y étaient fixés, qui passaient derrière sa tête, où on
les fermait à clef l'un à l'autre au moyen d'une petite clef pendue au bout d'une chaîne que
le Gardien des Portes portait autour du cou. Une fois les lunettes en place, Dorothy n'aurait
pas pu les ôter même si elle l'avait voulu ; mais bien sûr elle ne voulait pas être aveuglée
par l'éclat de la Cité d'Émeraude, et elle ne dit rien.}

\note[Choix de traduction 1]{\emph{a pair that would just fit Dorothy} : \emph{just} ne
signifie pas ici \og juste \fg{} au sens de \og seulement \fg, mais \og exactement \fg. Et
\emph{to fit} = \og aller à \fg{} en parlant d'un vêtement, non \og adapter \fg.}

\note[Choix de traduction 2]{\emph{could not take them off had she wished} : inversion à la
place de \emph{if she had wished}. Tournure littéraire ; le français rétablit \og même si elle
l'avait voulu \fg.}

\note[Choix de traduction 3]{\emph{a little key that was at the end of a chain the Guardian
wore around his neck} : deux relatives emboîtées, la seconde sans pronom. Le français ne peut
pas l'omettre — \og une chaîne \emph{que} le Gardien portait \fg.}

%% ===================================================================
\partiel{Lexique à mémoriser}
\begin{multicols}{2}\small
\begin{itemize}[label=--,itemsep=0.8mm,leftmargin=*]
  \item \textbf{spectacles} — des lunettes
  \item \textbf{a gate} — une porte (de ville, de jardin)
  \item \textbf{a shop} — un magasin
  \item \textbf{to sell} — vendre ; sold, sold
  \item \textbf{to buy} — acheter ; bought, bought
  \item \textbf{clothes} — les vêtements \textit{(toujours pluriel)}
  \item \textbf{a beard} — une barbe
  \item \textbf{a cart} — une charrette
  \item \textbf{to hide} — se cacher \textit{(fiche 7)}
  \item \textbf{to seem} — sembler, paraître
  \item \textbf{a key} — une clef
  \item \textbf{a neck} — un cou
\end{itemize}
\end{multicols}

\end{document}
